%%%%%%%%%%%%%%%%%%%%%%%%%%%%%%%%%%%%%%%%%%%%%%%%%%%%%%%%%%%%%%%%%%%%%%%%%%%%
\documentclass[a4j]{ujarticle}

\usepackage{docmas_v2}

%%%%%%%%%%%%%%%%%%%%%%%%%%%%%%%%%%%%%%%%%%%%%%%%%%%%%%%%%%%%%%%%%%%%%%%%%%%

% end of template style settings
%\usepackage{graphicx}%
\usepackage[dvipdfmx]{graphicx}
\usepackage{multirow}%
\usepackage{amsmath,amssymb,amsfonts}%
\usepackage{amsthm}%
\usepackage{mathrsfs}%
\usepackage[title]{appendix}%
\usepackage[dvipdfmx]{xcolor}%
\usepackage{textcomp}%
\usepackage{manyfoot}%
\usepackage{booktabs}%
\usepackage{algorithm}%
\usepackage{algorithmicx}%
\usepackage{algpseudocode}%
\usepackage{listings}%


\usepackage[normalem]{ulem}

\newcommand{\argmax}{\mathop{\rm arg\,max}\limits}
\newcommand{\argmin}{\mathop{\rm arg\,min}\limits}

\def\erase#1{\red{\sout{#1}}}
\def\subst#1#2{\dgreen{#1} \red{\erase{#2}}}

\def\blue#1{\textcolor[rgb]{0,0,1.0}{#1}}
\def\green#1{\textcolor[rgb]{0,0.6,0}{#1}}
\def\red#1{\textcolor[rgb]{1.0,0,0}{#1}}
\def\green#1{\textcolor[rgb]{0,0.6,0}{#1}}
\def\dgreen#1{\textcolor[rgb]{0.02,0.55,0.04}{#1}}
\def\gray#1{\textcolor[rgb]{0.5,0.5,0.7}{#1}}

\def\Dgreen{\color[rgb]{0.02,0.55,0.04}} \def\Blue{\color[rgb]{0,0,1}}
\def\Red{\color[rgb]{1,0,0}} \def\Black{\color[rgb]{0,0,0}}
\def\Gray{\color[rgb]{0.5,0.5,0.7}}
\def\colorreset{\color[rgb]{0,0,0}}
\def\subst#1#2{\dgreen{#1} \red{(削除 #2)\xspace}}

\usepackage{xspace}


% parameter macro definition

\def\AgentSet{\mathcal{A}}
\def\post{\rho} % was P_p
\def\BoundedConfidence{\epsilon}
\def\e{\varepsilon}
\def\feed{F}
\def\postOpVar{V_p}
\def\stubbornness{\mu}

% maxで8ページ



\begin{document}

% 和文タイトル
\title{SNSにおける中立的インフルエンサーの \\ 議論参加が引き起こす分極化加速メカニズムの解析}

% 英文タイトル
\etitle{
  %The Role of Moderate Influencers as Hubs \\ in Accelerating Opinion Polarization on Social Media
  } 

% 著者名: 
% ・各著者を\quad（全角空白）区切りで列挙
% ・著者名の直後に\afil{所属番号}を追加→所属番号を上付で出力（\textsuperscript{所属番号}と同じ）
% ・複数機関へ所属している場合は番号をカンマ区切りで列挙（下記著者２参照）
% ・Corresponding Authorについては所属の後に\thanksを続け，連絡先を記入
% ・英文著者はカンマ区切りで列挙

\author{武田朋也\afil{1}%
	\thanks{
  早稲田大学大学院 基幹理工学研究科 情報理工・情報通信専攻 菅原研究室 \newline%
  @E-mail: tomoya.leo9@fuji.waseda.jp}\quad%
  鳥海不二夫\afil{2}
  \thanks{東京大学大学院 工学系研究科 システム創成学専攻 \newline%
  @E-mail: tori@sys.t.u-tokyo.ac.jp}\quad%
  菅原俊治\afil{1}
  \thanks{早稲田大学大学院 基幹理工学研究科 情報理工・情報通信専攻 \newline%
  @E-mail: sugawara@waseda.jp}\quad%
  \\
	Tomoya Takeda\afil{1} \quad \quad Fujio Toriumi\afil{2} \quad \quad Toshiharu Sugawara\afil{1}}

\affiliation{
	\afil{1} 早稲田大学 理工学術院 基幹理工学研究科 情報理工・情報通信専攻\\ 
  \afil{1} Department of Computer Science and Communications Engineering，Waseda University\\
  \afil{2} 東京大学大学院 工学系研究科 システム創成学専攻\\ 
  \afil{2} Department of Systems Innovation，School of Engineering，The University of Tokyo\\
}

\abstract{
本研究では，対立する極端な意見を表明するインフルエンサーが存在する環境において，
中庸な意見を持つインフルエンサーの情報拡散活動の影響を調査する．
SNS上での意見分極に関する研究は幅広く行われてきたが，論争的な議論に
おいて異なる立場を持つ複数のインフルエンサーの
活動パターンが，議論全体のダイナミクスに与える影響を調査した研究は少ない．
本研究では，Bounded Confidenceモデルをはじめとした確立された社会科学モデル
をもとに，エージェントベースモデル(ABM)をスケールフリーな
ユーザネットワーク上に構築し，意見ダイナミクスのシミュレーションにより，
その影響を調査する．
特に，異なる特定の意見を持つインフルエンサーの投稿及びリポスト行動に着目して
活動パターンを制御することで意見への影響の比較実験を行う．
その結果，対立する極端な意見を表明するインフルエンサーが存在する環境において，
中庸な意見を持つインフルエンサーの情報拡散活動は，そのフォロワーのエコーチェンバー
の形成を助長し，結果的に全体の意見分極を加速させる可能性を示した．
これらの知見は，SNSをはじめとするオンラインソーシャルネットワーク上においての
ハブユーザの複雑な役割を明らかにすると同時に仲裁行動の難しさを推察させ，
将来のオンラインソーシャルネットワークプラットフォームの設計に対して
重要な示唆を与える．
}

\maketitle
\thispagestyle{empty}


\section{はじめに}\label{Intro}
% SNSの影響力と社会問題
ソーシャルネットワーキングサービス(SNS)は，現代社会では多くの人々
に利用され，その影響力は増し続けている．SNSの登場によって人々は従来ま
でとは異なる情報空間を経験し，SNSは高度情報化社会を代表するツー
ルとなった．現在ではX(旧Twitter)やFacebook，YouTube，TikTokなどをはじ
めとした様々なプラットフォームが存在し，そこでは日常的なコミュニケーショ
ンから，エンターテインメント，政治，社会問題に関する情報交換など多
くの情報が日々発信され，消費されている．SNSの普及により，ユーザは物理的な
制約を越えて円滑にコミュニケーションや情報収集が可能となり，世界中の不特定多
数に向けて容易な情報発信ができるようになった．新しい情報流通プラッ
トフォームとして多くの利点をもたらしたSNSだが，同時に様々な問題
が存在することも指摘されている．たとえば，ユーザの選好的な人間関係の選択により自
身を肯定するユーザとのやり取りが支配的になるエコーチェンバー現象やSNS
プラットフォームの推薦アルゴリズムにより，ユーザが好むコンテンツばか
りがタイムライン上に流れるフィルターバブル現象，多様な言論やコンテンツ
の淘汰，AIを組み込んだbotを用いたフェイクニュースの拡散や情報操作など
があり，これらの多様な問題が複雑に絡み合いながら顕在化してい
る~\cite{Fisher2015,Baccarella2018,Conover2021}．昨今は，デジタル民主
主義の考え方に基づき，オンラインソーシャルネットワークが市民による
熟議の場を実現する基盤技術として検討，導入されてい
る~\cite{Smith2021,Tseng2022}．今後も更なる変化が予想される情報化社会
を生きる者にとって，こうしたSNSの特徴や問題点を深く理解し，個
人や社会への影響を精査することを通して，社会の中でのSNSのあり方につい
て議論することは重要である．
\par


計算社会科学\cite{Lazer2009}分野に代表されるように，多くの研究がSNSの
抱える様々な問題を調査し，有用な知見を提供している．その中でも，SNS上
のユーザネットワークの主要な特徴であるスケールフリー性に着目し，インフ
ルエンサーがSNS上での議論に与える影響について調査した研究は存在す
る~\cite{Betts2023,Helfmann2023}．しかし意見の異なる複数のインフルエンサー
の活動パターンが意見分極に与える影響について調査した研究は少ない．
\par


本研究は，SNS環境をモデル化したマルチエージェントシミュレーションを用
いて，対立的な議論の中で特定の意見を持つインフルエンサーらの活動パター
ンを制御することで，インフルエンサーの活動パターンがSNS上での世論分極
に与える影響について調査した． その結果，対立する両極端な意見を持つイ
ンフルエンサーが存在する環境下では，中立的な意見を持つインフルエンサー
の情報拡散行動が，極端な意見を保持していたユーザらを更に閉鎖的なエコー
チェンバーへと陥れ，全体の意見分極を加速させる現象を確認した．この現象
は，中立的な立場を持つインフルエンサーのユーザネットワーク上でのハブ的
な役割と，対立する両極端なユーザ間での非需要的なコンテンツ消費という2
つの要素を用いて説明できる．本研究の発見は，SNSをはじめとするオンライ
ンソーシャルネットワークプラットフォームにおけるハブ機能の重要性や中立
的な仲裁行動の難しさを示唆しており，健全な言論空間を目指すプラットフォー
ムの設計に対する重要な問題提起となる．
\par

%

\section{関連研究}
社会システムにおける人々のつながりを複雑ネットワークとして捉え，それを
計算論的に表現することを通して社会現象に対する示唆を得ようと試みる研究
は，SNSが登場する以前から広く行われている．特に，社会を構成する個人を,
内在変数を持つ独立したエージェントとして定式化し，多様な人間関係をエー
ジェント間の相互作用として計算論的にモデル化するマルチエージェントシス
テム(MAS)を用いた社会シミュレーションはその代表的手法であり，単純なモ
デル設計でも複雑なダイナミクスを表現できるため，現在でも多くの研究で利
用されている．たとえばLiggettら\cite{Liggett1975}は，グラフネットワー
ク上のノードに二値変数を持たせ，それをノードが接続する隣接ノードが持つ
変数の値に影響される形で確率的に変化させることを通して，系全体での状態
や収束について数学的に解析できるモデル(Voterモデル)を提案した． Voter
モデルでは，グラフ上のノードを，ある論点に対して賛成または反対という二
値の意見を持つエージェント，ノード間のリンクをエージェント間の人間関係
と捉えることで，系全体の変化は世論のダイナミクスを表現していると考えら
れる．このように世論のダイナミクスについて調査する研究トピックは,
Opinion Dynamicsと呼ばれ，Voterモデルはその基礎的な提案である．
\par


この後にMASを用いたOpinion Dynamicsに関する研究が行われ，社会科学的な
視点を導入したモデル構築がなされるようになった． Voterモデルにおいては,
あるエージェントの意見は隣接エージェントの意見に影響されて変化するが,
その相互作用 (Social Influence)について注目して拡張した研究として
\cite{DeGroot1974,Friedkin1990,Hegselmann2002}がある． DeGrootら
\cite{DeGroot1974}は，人間は自分の意見を保持しつつも他者の意見を取り入
れることをモデル化し，ある瞬間のエージェントの意見を，前の瞬間の隣接エー
ジェントの意見の重み付き平均として算出するモデルを考案した．さらに
Hegselmannら\cite{Hegselmann2002}はDeGrootモデルを拡張し，他者との意
見の差が大きすぎる場合は影響を受けないという仮説のもとで信頼閾値
(Bounded-Confidence)と呼ばれる閾値を導入したモデルを提案した．
\par


現在，SNSの登場により従来とは異なる人間関係が生まれ，それ
らを対象にした研究も盛んに行われている．笹原ら\cite{Sasahara2021}は,
SNS環境を想定したシミュレーションモデルを用いることで，他者の意見に影
響を受けるというユーザの社会的特徴と，unfollowやblockというSNS特有の人
間関係の拒絶行動の組み合わせが分極を生じさせることを明らかにし，そのシ
ミュレーション結果と意見分極化下のSNSの実際のログデータ
\cite{Conover2021}とで整合性が取れることを確認した．このように情報社会
の発展につれ，Opinion Dynamicsをはじめとした社会シミュレーショ
ン研究は，人々の行動ログに関するビッグデータを用いることで定量的な検証
が可能になった．複雑な社会現象を社会科学的な視点を用いて社会シミュレー
ションやビッグデータ解析などのアプローチで研究する学際的領域は計算社会
科学 \cite{Lazer2009}と呼ばれ，分野横断的な知見を織り込んだ研究領域と
して体系立てられている．
\par


SNSに関する研究の中でも，特にインフルエンサーの行動や役割について調査
した研究は多い． Kempeら\cite{Kempe2003}やKumarら\cite{Kumar2022}は，
ユーザネットワークの中で，どのようなユーザ集合から情報を発信すれば情報
拡散が最大化されるかという影響最大化問題に対する現実的な計算法を考案し
た．Helfmannら\cite{Helfmann2023}は，メディアとインフルエンサーと一般
ユーザを区別したモデルを用いることで，インフルエンサーが極端な言動に傾
くことがフォロワーを増やす効率的な戦略であることを示した．Bettsら
\cite{Betts2023}は，インフルエンサーと一般的なユーザとを明確に分離した
SNSシミュレーションモデルを構築し，その結果から中立的で広く支持を集め
るカリスマ的なインフルエンサーによる仲介であれば，意見分極を抑制する効
果があると主張した．これらの研究をはじめとして，SNS上のインフルエンサー
に着目した研究は多いが，複数の立場のインフルエンサーの様々な活動パター
ンが議論全体の意見分極に与える影響について調査した研究は少ない．
\par


\section{提案モデル}
\subsection{Agent-Based Model}
本節では，SNSプラットフォームにおける意見ダイナミクスをモデル化した
Agent-Based Model (ABM) を導入する．$G=(\AgentSet,E)$ をユーザネットワー
ク構造を表す有向グラフとする．ここで，$A=\{1,2,\ldots,N\}$ は $N$ 人の
エージェント（あるいはユーザ）の集合を表し，$E \subseteq \AgentSet
\times \AgentSet$ は「フォロー」関係を表す有向エッジの集合である．有向
エッジ $(i,j)\in E$ は，エージェント $i$ がエージェント $j$ をフォロー
していることを意味する．隣接行列 $\mathbf{A}$ を，$A_{ij}=1$ if
$(i,j)\in E$，and $A_{ij}=0$ otherwiseと定義する．

  
\begin{table}[t]
  \tabcolsep=5pt
  \centering
  \caption{エージェントが持つパラメータ}\label{tab:agentparameters}
  \begin{tabular}{llcl}
    \toprule
    パラメータ & 記号 & 範囲 & 初期値 \\
    \midrule
    意見 & $o_i(t)$ & $[-1.0，1.0]$ & $\mathcal{N}(0.0,\,0.6)$ \\
    投稿確率 & $\post_i(t)$ & $[0.05，1.0]$ & $0.1$ \\
    信頼閾値 & $\BoundedConfidence_i(t)$ & $[0.2，1.0]$ & $1.0$\\
    \bottomrule
  \end{tabular}
\end{table}

非負整数$t \in \{0,1,2,\ldots\}$で表される離散時間を導入する．
各エージェント $i \in A$ は，表\ref{tab:agentparameters}に示す内部パラ
メータによって特徴づける．意見 $O_i(t) \in [-1,1]$ は，時刻 $t$ におけ
るエージェント $i$ の特定の論点に対する正規化された立場を表す．境界値
$-1$ と $1$ は互いに対立する最も極端な意見を表し，$0$ は中立的あるいは
穏健な態度を表す．初期意見 $O_i(0)$ はエージェントの「潜在的意見
（intrinsic opinion）」とし，その基礎的信念を表すもので，他
エージェントとの相互作用とは関係なくエージェント固有の値として自身を特
徴づけ続ける．
\par


投稿確率 $\post_i(t)$ は，エージェント $i$ がある時刻$t$において
投稿を行う確率である．このパラメータは，社会シミュレーション研究でよく
用いられる「強度（intensity）」の概念に対応する．本モデルの重要な点は，投稿強度 $\post_i(t)$ が意見の極端さ $|O_i(t)|$ とは独立に
定義されることである．すなわち，活動強度は過激化の度合いに必ずしも比
例しないことを想定している．この独立性により，「穏健な意見を持つユーザ
が高頻度で投稿する」といった多様な行動パターンを捉えることが可能とな
る．
\par


第\ref{Intro}節で述べたように，Bounded Confidenceモデルは，
個人は自分の意見から大きく乖離した意見には影響されにくいという社会心理
学的仮説を表現している．
本モデルでは静的な閾値ではなく，動的な信頼閾値パラメータ
$\BoundedConfidence_i(t)$ を導入する．
これにより，プラットフォーム上での相互作用を通じて，
ユーザが対立する意見に対してどの程度開かれているかという「許容度」が時
間とともに適応的に変化する様子をモデル化できる．
これらのパラメータの進化は，次節で説明する社会的更新メカニズムによって
規定される．
\par

% FJについてはここで述べる．degrootの拡張として載せる．



\subsection{社会的適応のモデル化}
一般に，人間は自身の内部状態を周囲の環境に適応させ
る~\cite{NoelleNeumann1974,Cialdini2004}．我々は，これまでの研究で確立
されてきた計算理論に基づき，これらのメカニズムをモデル化する．エージェ
ントは一定確率 $P_a$ でプラットフォームにアクセスするたびに，以下のメ
カニズムを通じて内部パラメータを更新する．ここで，エージェントの信頼閾
値を用いて，\emph{受容的な投稿 (comfortable posts)} と \emph{非受容的
  な投稿 (uncomfortable posts)} を定義する．具体的には$|O_i - p_j|
\leq \BoundedConfidence_i$ のとき投稿 $p_j$ は受容的であり，$|O_i -
p_j| > \BoundedConfidence_i$ のとき非受容的とする．


\begin{description}
\item \textbf{Social Influence:} 本モデルでは，ユーザの意見の更新に関
  して，ユーザがタイムラインで閲覧した投稿の意見と自身の潜在的な意見を
  平均化することで意見を更新するというFriedkinらによって定式化されたモ
  デル~\cite{Friedkin1990}を採用する．エージェント$i$はタイムステップ
  $t$において，自身のタイムライン$F_i(t)$内のすべての投稿が持つ意見の
  値と，エージェント自身の潜在的意見$O_i(0)$の重み付き平均を算出するこ
  とで意見を更新する．具体的には，以下のように定式化する．
  %
  \begin{equation}\label{FJmodel}
   O_i(t) = \stubbornness \cdot O_i(0) + (1 - \stubbornness) \cdot \frac{\sum_{p_j \in F_i(t)} O(p_j)}{|F_i(t)|}
  \end{equation}
  %
  ここで $\stubbornness$ はエージェントの\emph{頑固さ (stubbornness)}
  を表すパラメータであり，他者に影響されることに対して，エージェントが
  自身の潜在的意見 $O_i(0)$ にどの程度固執するかを示す．今回の実験では，
  $\stubbornness$は全エージェントに対して一定の値で固定する．  なお，
  後に説明するが，エージェントのタイムライン内の投稿はすべて，そのエー
  ジェントがフォローしているユーザーが投稿またはリポストしたものである．
  したがって，エージェントはその作成者をアンフォローまたはブロックしな
  い限り，非受容的な投稿を閲覧する可能性がある．

\item \textbf{Social Reinforcement:}\label{SocReinfo} ユーザはタイムラ
  イン上に自身と似た意見を主張する投稿が多いと，自身の投稿行動を活発化
  させる．本モデルでは，特に沈黙の螺
  旋~\cite{NoelleNeumann1974,Sohn2015}に関するエージェントシミュレーショ
  ン研究で用いられる閾値の概念を導入する．エージェントは，タイムライン
  内の受容的な投稿の割合に基づいて，全エージェントに共通な閾値
  $\epsilon_s$と比較することで自身の投稿確率を増減させる．形式的には，
  $K_i(t)=\{p_j \in F_i(t) \mid |O_i(t) - O(p_j)| \leq
  \BoundedConfidence_i\}$ としたとき，$C_i(t)= |K_i(t)|/|F_i(t)|$ をタ
  イムステップ $t$ におけるエージェント $i$ のタイムライン内の受容的な
  投稿の割合として，更新ルールは以下の通りである．
  %
  \begin{align}
    \begin{aligned}
      \text{If } C_i(t) & \geq \epsilon_s，\\[-2pt]
      & \post_i = \min(\post_i \cdot (1 + \delta_{\post})，1.0)\\
    \end{aligned}\\
    \begin{aligned}
      \text{If } C_i(t) & \leq 1-\epsilon_s，\\[-2pt]
      & \post_i = \max(\post_i \cdot (1 - \delta_{\post})，0.05)
    \end{aligned}
\end{align}
  %
  ここで $\epsilon_s$ は意見の普及度 (opinion prevalence) に対する閾値
  を表し，$\delta_{\post} > 0$ は強度の調整率を制御するパラメータであ
  り，すべてのエージェントに対して共通の値とする．

\item \textbf{Confidence Decline:}\label{ConfDec} 本モデルでは信頼閾値
  $\BoundedConfidence$の更新について，肯定的な経験よりも否定的な経験の
  方が大きい心理的重みを持つ~\cite{Vallone1985,Baumeister2001}という非
  対称な人の性質を反映させる．エージェントは，以下のように，タイムライン内の非受容的
  な投稿を読むたびに，自身の信頼閾値を減少させる．
  %
  \begin{equation}
    \begin{aligned}
    \text{If } & |O_i(t) - p_j| > \BoundedConfidence_i，\\[-2pt] 
    & \BoundedConfidence_i = \max(\BoundedConfidence_i \cdot (1 - \delta_{\BoundedConfidence})，0.2)
    \end{aligned}
  \end{equation}
  ここで $\delta_{\BoundedConfidence} > 0$ は，非受容的な周囲の環境に
  対する信頼の減衰率を表すパラメータであり，全エージェントで共通の値とする．
\end{description}


\subsection{SNS上でのユーザ行動のモデル化}
SNS環境でのユーザによる意見形成ダイナミクスを再現するため，その特徴的な行動として，投稿，リポスト(いいね)，フォロー，アンフォロー(ブロック)をモデル化する．


\begin{description}
\item \textbf{投稿:} エージェント$i$はプラットフォームにアクセスするた
  びに，確率 $\post_i$ で自身の現在の意見$O_i(t)$を反映した投稿を作成
  する．エージェント$i$が投稿する，その投稿は$i$をフォローしているエー
  ジェントのタイムラインに追加される．
\item \textbf{リポスト/いいね:} エージェント $i$ はタイムライン内の$j$
  からの投稿$p_j$を閲覧した際，それが受容的（すなわち $|O_i - O_{p_j}|
    \leq \BoundedConfidence_i$）である場合，確率 $P_r$ でリポストする．
  エージェント$i$によってリポストされた投稿は，$i$のフォロワーのタイム
  ラインに追加される．投稿はこのリポスト機能を通じてSNSネットワーク全
  体に拡散される．なお，実際のSNSサービスにおける「いいね」機能は，こ
  こで述べたものと同様の拡散効果を持つため，本モデルでは「いいね」と
  「リポスト」を区別しない．

\item \textbf{フォロー:} 同様に，エージェント$i$が自身のタイムライン内の
  $j$からの投稿$p_j$を閲覧した際に$p_j$が受容的であり，かつ$p_j$の作成者$j$との
  間にフォロー関係が現在存在しない場合，確率 $P_f$ でその投稿の作成者$j$
  をフォローする．
  
\item \textbf{アンフォロー/ブロック:} 逆に，エージェント$i$が投稿$p_j$を
  非受容的（すなわち $|O_i - O_{p_j}| > \BoundedConfidence_i$）と捉え
  た場合，その投稿の作成者$j$との間にフォロー関係が存在すれば確率 $P_u$
  で$j$をアンフォローし，関係が存在しない場合は，同じく確率 $P_u$
  で作成者$j$をブロックし，どちらの場合でも$j$からの以降のすべての
  投稿を拒絶する．
\end{description}


%どのようにインフルエンサーを制御するか
\subsection{インフルエンサーの活動制御}
本研究の主要な着眼点であるインフルエンサーの活動パターンを，モ
デル内において以下のように制御する．初期ネットワーク生成後，各エー
ジェントの流入次数（フォロワー数）に基づき，上位3つのノードか
らなる集合を $\AgentSet_{T} \subset \AgentSet$ ($|\AgentSet_{T}| = 3$)
と定義する．集合 $\AgentSet_{T}$ に属するエージェントに対して特定
の意見値および行動確率を割り当てることで，ネットワーク構造上のハブエー
ジェントの行動パターンが全体の意見ダイナミクスに与える影響を調査する．
具体的には，$\AgentSet_{T}$に属する3エージェントの初期意見(潜在的意見)
を(-1.0，0.0，1.0)の3つの代表値から実験試行ごとにランダムに割り当て，
これら3エージェントの投稿確率及び，リポスト確率を制御することで，複数
のオピニオンリーダーの様々な行動パターンをシミュレーションできる．


\section{実験と考察}
\subsection{実験設定}
SNSのユーザネットワーク環境$G = (\AgentSet，E)$の初期構造をCNNRネット
ワーク（Connecting Nearest Neighbor with Random linkage: CNNR）
\cite{Yuta2007}を採用する．CNNRモデルは，CNNモデル~\cite{Vazquez2003}
を拡張したネットワーク生成アルゴリズムである． CNNモデルは社会ネットワー
クの3つの主要な特性であるスケールフリー性，クラスタ性，スモールワール
ド性を備えることで知られる．これにランダムな接続を加えることで，より現
実的な社会ネットワークを構築できると考えられる．なお，本実験では，潜在
エッジから実エッジを生成する確率を $u = 0.3$，ランダムな接続確率を $r
= 0.01$ とした．
\par


表\ref{tab:expparameters}に実験で用いた各実験パラメータの値を，表
\ref{tab:networkparameters}にCNNRネットワークの生成アルゴリズムによ
り生成された初期ネットワークの各種統計量の平均値を列挙する．
実験では，これまで説明してきた実験環境で，十分にダイナミクスが収束する40000ステップを最終ステップとして実験を行った．
\par

\begin{table}[t]\tabcolsep=5.0pt
  \begin{minipage}[t]{0.2\textwidth}
   \centering
   \caption{実験パラメータ}\label{tab:expparameters}
   \begin{tabular}{llcl}
     \toprule
     パラメータ & 値 \\
     \midrule
     $P_a$ & 0.1 \\
     $P_r$ & 0.4 \\
     $P_f$ & 0.1 \\
     $P_u$ & 0.1 \\
     $\stubbornness$ & 0.8\\
     $\epsilon_s$ & 0.5\\
     $\delta_p$ & 0.1 \\ 
     $\delta_\BoundedConfidence$ & 0.01 \\
     \bottomrule
   \end{tabular}
  \end{minipage}
  \hfill
  \begin{minipage}[t]{0.3\textwidth}
    \centering
    \caption{ネットワーク統計量}\label{tab:networkparameters}
    \begin{tabular}{llcl}
      \toprule
      パラメータ & 値 \\
      \midrule
      エージェント数 & 1000 \\
      平均次数 & 1.42 \\
      平均経路長 & 3.23 \\
      直径 & 11.0 \\ 
      べき乗指数 & 2.46 \\
      クラスタ係数 & 0.233 \\
      \bottomrule
    \end{tabular}
  \end{minipage}
  \end{table}


\subsection{オピニオンリーダーの活動パターンと意見分極}
本実験では，$\AgentSet_{T}$ に属し，$(-1，0，1)$のいずれかの意見を持つ
3つのエージェント（以下，オピニオンリーダー）の行動パターンが意見分極
に与える影響を明らかにするため，各オピニオンリーダーの活動状態を制御し
た比較実験を行った．具体的には，$\AgentSet_{T}$ に割り当てられた意見
$O \in \{-1.0，0.0，1.0\}$ を持つ各オピニオンリーダーに対し，投稿確率
$\post$ およびリポスト確率 $P_{r}$ を以下の2水準のいずれかに設定した．

\begin{itemize}
\item \textbf{Active (A):} 投稿確率およびリポスト確率をモデルの規定
  通り（0に固定せず）に設定し，議論に参加する状態．
\item \textbf{Inactive (I):} 投稿確率およびリポスト確率を常に $0$ に
  固定し，情報の拡散・発信を一切行わない状態（単なるフォロワーとして
    振る舞う）． 
\end{itemize}

$\AgentSet_{T}$ は異なる意見（$-1.0，0.0，1.0$）を持つ3体のエージェン
トで構成されるため，各オピニオンリーダーの活動状態（A または I）の組み
合わせは合計 $2^3 = 8$ 通り存在する．本稿では，各実験条件におけるオピ
ニオンリーダーの活動パターンを，意見 $-1.0，0.0，1.0$ の順に並べた組
$(s_{-1}，s_{0}，s_{+1})$ で表記する．ここで $s \in \{A，I\}$ である．
例えば，条件 $(A，I，I)$ は，意見 $-1.0$ を持つオピニオンリーダーのみ
が活動 (Active) し，中庸 ($0.0$) および対立意見 ($1.0$) を持つオピニオ
ンリーダーは活動しない (Inactive) 状況を表す．これら8通りの組み合わせ
について，異なる乱数シードを用いた10回の独立した試行を行い，そ
の平均的な挙動を分析した．Opinion Dynamicsの分野では意見分極を評価する
指標として，しばしばユーザが持つ意見の分散値が用いられる．そこで本研究
においても，意見分極の程度を測る定量的指標として，ある期間に投稿された
すべての投稿に反映されている投稿作成者の意見の分散値を採用し，これを
$\postOpVar$と定義する．
\par


図\ref{fig:var_results} は，全8通りの実験条件について，各10回の試行に
おける40000ステップ直前の最終100ステップで投稿された投稿の意見の分散値
$\postOpVar$ の平均値を示している．この結果から，意見分極が
相対的に激しい状態に至ったのは，実験設定 $(A，A，A)$や$(A，I，A)$など
の両極端な意見を持つオピニオンリーダーが双方とも活動している場合であり，
逆に意見分極の度合いが小さかったのは$(I，A，I)$や$(I，I，I)$などの両極
端のオピニオンリーダーが活動していない実験設定であった．また，注目すべ
きは実験設定$(A，I，A)$に比べて，中立的な意見を持つオピニオンリーダー
が活動する$(A，A，A)$の方が分極が悪化しているという結果である．これら
の現象に対する解釈を含め，以下では極端なオピニオンリーダーと中立的オピ
ニオンリーダーのそれぞれに着目し，実験設定ごとの結果を詳細に解釈する．
\par

\begin{figure}[t]
  \centering
  \includegraphics[width=1.0\linewidth]{fig/var_results.pdf}
  \caption{条件ごとの終状態における意見分極度}
  \label{fig:var_results}
\end{figure}

\subsubsection{極端なオピニオンリーダーによる分極化の促進}
ここでは，$-1.0$もしくは$+1.0$という極端な意見を持つオピニオンリーダー
(以下，極端なオピニオンリーダー)が意見分極に与える影響について考察する．
$(I,I,I)$ や $(I，A，I)$ などの極端なオピニオンリーダーが活動しない
環境と比較して，$(A，I，I)$ や $(A，A，I)$ のように極端なオピニオンリー
ダーが活動する環境の方が $\postOpVar$ の値が大きい．このこ
とから，$-1.0$ または $+1.0$ の極端なオピニオンリーダーの活動は，それ
らが存在しない環境に比べて意見分極を増大させる効果があると言える．
\par


この現象は，極端なオピニオンリーダーの活動が，彼らに似た意見を持つエー
ジェントのエコーチェンバー形成を促進するというメカニズムで説明で
きる．極端な意見を持つ中・小規模のエージェントは，対立する意見
を主張する投稿に接しやすく，特に自身の信頼閾値を超える非受容的な
投稿を閲覧する可能性が高い．第\ref{ConfDec}節で定義したConfidence
Declineメカニズムにより，このような非受容的な投稿への接触は連鎖的な信
頼閾値の減少を引き起こす．信頼閾値の減少は，自身と意見が乖離するエージェ
ントをアンフォローする行動を促し，結果として自身の意見に近いエージェン
トのみとの繋がりを強化する．こうして形成されていく閉鎖的な情報空間（エ
  コーチェンバー）内では，タイムラインにおいて自分と似た意見の投稿が支
配的となっていく．本モデルのSocial Reinforcementメカニズム
（第\ref{SocReinfo}節）に基づけば，タイムライン内の投稿の中で意見齟齬が
自身の信頼閾値より小さい受容的投稿の割合が高いほど投稿確率
$\post$ は上昇するため，エコーチェンバーに陥った極端なエージェントほど，
その発信活動を活発化させる．これが極端なオピニオンリーダーの活動がエ
コーチェンバーの形成を助長させて全体の意見分極を悪化させる基本的なメカ
ニズムである．
\par


\subsubsection{中立的オピニオンリーダーが招く分極の悪化}
次に，中立的な意見を持つインフルエンサー(以下，中立的オピニオンリー
ダー)が意見分極に与える影響について考察する．図~\ref{fig:var_results}
において特に注目すべき点は，中立的オピニオンリーダーが活動する $(A，A,
A)$ の結果が，中立的オピニオンリーダーが活動しない $(A，I，A)$ よりも
大きな $\postOpVar$ を記録していることである．これは，対立的な極端なオ
ピニオンリーダーが存在する環境下において，中立的なオピニオンリーダーの
活動が，意見分極を悪化させることを示唆している．
\par


ここで，極端なオピニオンリーダーが活動しない条件下，すなわち $(I，I,
I)$ と $(I，A，I)$ を比較すると，両者の $\postOpVar$ には大きな差が見
られない．したがって，中立的オピニオンリーダーの情報拡散行動そのものが
必ずしも分極を悪化させるわけではなく，「対立する両極端なハブが存在する」と
いう条件と組み合わさると分極を助長すると言える．
\par


この結果は，中立的オピニオンリーダーの持つ「ネットワークハブとしての機
  能」と「エコーチェンバー形成過程」により説明できる．前述の通り，極端
な意見を持つエージェントは信頼閾値の連鎖的な減少に陥りやすく，対立意見
を持つエージェントをアンフォローして自身の情報空間から排除する傾向にあ
る．対照的に，中立的オピニオンリーダーを含む中庸な意見を持つエージェン
トは，多様な意見にバランスよく触れるため，片方に偏った非受容的接触が相
対的に少なく閾値の低下が緩やかであるため，$-1.0$側と$+1.0$側
の双方の意見のエージェントをフォローし続ける可能性が高い．結果とし
て，中立的オピニオンリーダーは双方の陣営からの投稿を，自身のリポスト行
動を通じて，それらを再拡散する橋渡し(ハブ)の役割を果たすことにな
る．
\par


このハブ機能により，ネットワーク内には両極端な意見が広く流通する．しか
し，すでに信頼閾値が小さく，排他的になりつつある極端なエージェ
ントにとって，中立的オピニオンリーダーを経由して流れてくる対立陣営の意
見は，非受容的でアンフォローすべき対象として認識される可能性が高い．結
果として，中立的オピニオンリーダーが活発に多様な意見を媒介すればするほ
ど，中・小規模の極端なエージェントが対立意見に接触する頻度は高まり，信
頼閾値の減少とアンフォロー行動が加速し，更に閉鎖的なエコーチェンバーに
陥る．以上のメカニズムにより，対立する両極端なオピニオンリーダー
が存在する環境において，中立的オピニオンリーダーの活動はそのハブ的役割
によって，極端なエージェントのエコーチェンバー形成を促進し，結果的に意
見分極を深刻化させると考えられる．
\par


現実世界に対応づけて考えると，これまで議論してきた結果は，SNS上におけ
る，ある対立的な議論に対するインフルエンサーの仲裁的な議論への介入や興
味本位での議論への介入が意図せずして更なる意見分極の悪化を招く可能性を
示唆する．この逆説的な結果を招く主要な要因は，これまでの説明のように，中
立的オピニオンリーダーの中立的な立場であるからこそ維持される多様なフォ
ロワー属性である．つまり，
中立的オピニオンリーダーのネットワークトポロジー上での立ち位置が重要で
あり，彼らの善意の有無や振る舞いの仕方に関係なく，彼らは意見分極の活性
化に寄与するハブとなりうる．実際のSNS上での実験においても，ユーザを対
立する意見に接触させることが政治的な分極を悪化させるという研究結
果\cite{Bail2018}が存在し，本実験の結果とも併せて，オンラインソーシャ
ルネットワーク環境下では，意見対立の融和のために，対立する意見を持つユー
ザ同士を単純に接触させるだけでは効果が見通しにくいことが強く示唆される．
\par


\section{結論}
本研究では，代表的な社会科学モデルを基にしたSNS上での議論に関するマル
チエージェントシミュレーションモデルを構築し，スケールフリーなユーザネッ
トワーク上において，特に，影響力の大きいオピニオンリーダーの活動パター
ンに着目して意見分極に関する調査を行った．その結果，極端な意見を持った
オピニオンリーダーが議論に参加することが，ユーザのエコーチェンバー形成
を助長することを通して，全体の意見分極を悪化させることを確認した．さら
に，そのような極端なオピニオンリーダーが議論に参加している環境において，
中立的な意見を持つインフルエンサーの議論への参加は，そのネットワーク上
でのハブ機能によって，ユーザのより閉鎖的なエコーチェンバー形成を促す効
果があり，結果的に意見分極を更に悪化させることを明らかにした．この結果
は，SNSをはじめとしたオンラインソーシャルネットワークプラットフォーム
におけるハブ機能の重要性と，対立する議論の仲介行動の難しさを示唆する．
\par


今回用いたSNSでの意見ダイナミクスに関するシミュレーションモデルは，複
雑な要素を排してオピニオンリーダーの活動パターンに注目した議論をするた
めに，実際のSNS空間を大きく単純化しており，SNSに特徴的
な推薦システムや，より複雑なユーザの社会的特徴などを反映していないとい
う限界が存在する．今後の研究では，それら本モデルで考慮されていない特徴
による影響を鑑みることや，実際のSNSのデータとの比較を通した頑健性に関
する検証などが期待される．
\par



\section*{謝辞}
本研究は，JST ERATO JPMJER2502および科学研究費補助金(25K03188，22H03906)の支援を受けたものである．


%%%%%%%%%%%%%%%%%%%%%%%%%%%%%%%%%%%%%%

%\begin{thebibliography}{99}
%\small
% \bibitem{example} Example，I.: An Example of References，in {\it
%	 Proc．Int．Conference on Examples of Papers}，pp.~1--10 (2011)
%\end{thebibliography}

\small
\bibliographystyle{plain}
\bibliography{refer.bib}

\end{document}
